\section{Conclusion}

Topographic feature edge weights translate the intuitive concept of ``communities separated by mountains belong in different districts'' into a verifiable, non-partisan, algorithmically precise signal for the METIS graph partitioner. The key finding is that ridge-crossing penalties eliminate cross-watershed districts in Appalachian states at a modest compactness cost (4.1\% for NC mountain districts), while flat-state graceful degradation ensures the signal does not produce spurious effects in states without meaningful elevation variation.

The topographic weight modifier joins the M.1--M.7 signal library as a tool for redistricting commissions in mountain states — particularly North Carolina, Virginia, West Virginia, Tennessee, Pennsylvania, Colorado, Montana, and Idaho — where community-of-interest preservation along topographic lines is both legally supported and factually meaningful.
